% Round 18. ci_global branch. Activation requires independently reviewed PG18.
\section{Complete heights in a product ambient variety beyond complete intersections}
\label{r18pa:section}

This application has a different generic-fiber geometry from the
complete-intersection families.  Its ambient variety is
$\mathbf P^1\times\mathbf P^4$ and its smooth divisors have
second Betti number two.  They therefore cannot be isomorphic to
projective complete intersections of dimension four.  The new
local Green input is exactly the fourfold product-smoothing
theorem PG18; its independent activation is a dependency of
the theorem below.  All the remaining geometric, arithmetic
and global-height steps are proved here or use the already
frozen interfaces indicated explicitly.

\subsection{The native product, sections and geometric connectedness}

Put
\[
 \mathcal Y=\mathbf P^1_{\mathbf Z}\times\mathbf P^4_{\mathbf Z},
 \qquad\mathcal H=\mathcal O_{\mathcal Y}(1,1),
 \qquad h=c_1(\mathcal H).
 \tag{PA18.1}\label{r18pa:ambient}
\]
It is smooth and projective over $\mathbf Z$ with geometrically
integral fibers.  Write $x,y$ for the two native hyperplane
classes, so
\[
 x^2=0,\quad y^5=0,\quad\deg(xy^4)=1,
                  \qquad\deg(h^5)=5.
 \tag{PA18.2}\label{r18pa:degree}
\]
Fix an integer $D\ge5$ and let $V_D$ be the space of
bihomogeneous forms of bidegree $(D,D)$.

The Segre embedding embeds $\mathcal Y$ in $\mathbf P^9$,
with $\mathcal H$ the restricted hyperplane bundle.
Every bidegree-$(D,D)$ monomial is a product of $D$
bidegree-$(1,1)$ monomials: pair its $D$ factors in
the first set of variables with its $D$ factors in
the second.  Thus ambient degree-$D$ polynomials
surject onto $V_D$ over every field.
The one-point first and second jets and the two-point
first jets used in CG18/CE18 are therefore all supplied
by this actual section space for $D\ge3$.

Here the product ambient is not a projective complete
intersection.  The following computation supplies
the connectedness input directly, without applying
a complete-intersection hypothesis to it.  For
every positive $t$, the Kunneth formula gives
\[
        H^j(Y_{\bar k},\mathcal O_Y(-t,-t))=0
                                      \quad(0\le j<5).
 \tag{PA18.3}\label{r18pa:negative-vanishing}
\]
Indeed $H^0(\mathbf P^1,\mathcal O(-t))$ and
$H^0(\mathbf P^4,\mathcal O(-t))$ vanish;
on $\mathbf P^1$ only degrees zero and one
can occur, and on $\mathbf P^4$ only degrees
zero and four can occur.  The only possible
nonzero tensor product is therefore in total
degree five.  This calculation holds in every
characteristic.

Suppose $r\le4$ sections of positive powers
of $\mathcal H$ cut properly on $Y$, and
let $Z$ be their scheme-theoretic intersection.
The Koszul resolution, together with
\eqref{r18pa:negative-vanishing}, gives
\[
                         H^0(Z_{\bar k},\mathcal O_Z)=\bar k.
 \tag{PA18.4}\label{r18pa:connected}
\]
In fact the only additional cohomology in the
resolution would have total degree $5-j\ge1$
for its $j$th term with $j\le r\le4$;
the sole degree-zero contribution is the
constant class from $\mathcal O_Y$.
Thus every such smooth positive-dimensional
intersection is geometrically integral.

A divisor of class $D h$ with at most isolated
singularities is also geometrically integral.
It is a Cartier divisor on a smooth variety,
hence $S_2$; isolated singularities in dimension
four give $R_1$, so it is normal.
Equation~\eqref{r18pa:connected} for $r=1$
makes it connected, hence integral.
This is the precise replacement for the
projective complete-intersection integrality
lemma in the arithmetic construction below.

\subsection{The local PG18 input and deep monodromy}

At a finite valued field $K$ with ring $R$, take
$D$ actual sections $f_i$ of $\mathcal H_R$ and
an actual section $Q$ of $\mathcal H_R^D$.
Require that their special reductions have the
following properties:
\begin{itemize}
\item the $D$ divisors form a geometric SNC divisor;
      every positive-dimensional intersection is
      smooth of the expected dimension, every
      fivefold intersection is finite etale, and
      no sixfold intersection is nonempty;
\item $\bar Q$ meets every positive-dimensional
      factor stratum transversely and avoids
      every zero-dimensional factor stratum;
\item the generic divisor
      $X=V(\prod_i f_i+\pi Q)$ is smooth and
      geometrically integral.
\end{itemize}
All positive-dimensional factor strata and the
positive-dimensional centers of the ordered
model are geometrically integral by
\eqref{r18pa:connected}.  The quadruple and
quintuple intersections, in particular, are
retained as actual schemes.
The same Koszul computation with one of the
sections taken to be $Q$ proves geometric
integrality of every positive-dimensional
$H_I\cap V(\bar Q)$ required by PG18.

The precise local input used here is PG18's
full ordinary-actual G theorem for this model:
on a fourfold generic fiber the edge divisor
classes are generated by the two ambient
classes and the specified exceptional classes,
PG18's primitive triple argument kills $U_2$,
and its curve-comparison argument closes the
compatible-representative gate.  Its ambient
cohomological hypothesis holds because
$H^2(Y_{\bar k},\mathbf Q_\ell(1))$ is the
two-dimensional space of the actual native
divisor classes $x,y$.  The same property
holds over every finite residue extension;
PG18 must therefore be applied afresh to
all cycles after each unramified valued
extension.  None of these new cycles is
assumed to descend as a pullback.

For clarity, the top monodromy count does
not require a projective complete-intersection
ambient.  Every connected positive-dimensional
stratum contributes one copy of the constant
cohomology group.  Each fivefold intersection
has exactly $\deg(h^5)=5$ geometric points.
The ordered resolution does not change these
points or the curves immediately above them,
because $Q$ avoids the zero-dimensional
original strata.  Its weight-zero restriction
complex is the simplex complex through degree
three, with each top subset repeated five
times.  Its last differential has rank
$\binom{D-1}{4}$, so
\[
 \operatorname{rank}N^4
       =5\binom D5-\binom{D-1}{4}
       =(D-1)\binom{D-1}{4}>0.
 \tag{PA18.5}\label{r18pa:monodromy}
\]
The equality with the actual rank follows
by the HM18/NG18 extreme restriction--Gysin
argument: the opposite extreme matrix is
the transpose under the positive native
degree pairing on the individual geometric
point summands, and $N^4$ is the identity
on the common summands.  It gives an
isomorphism of the opposite extreme kernel
with the restriction cokernel.  Purity gives
$E_2$ degeneration, and the weights of
$H^4$ lie in $[0,8]$, so the extreme rank
is the actual fourth-monodromy rank.
No general weight-monodromy assertion is
invoked.  Closed-point residue degrees are
retained as the sizes of their geometric
point orbits.

\subsection{An arithmetic family at every finite place}

\begin{theorem}[Complete heights beyond projective complete intersections]
\label{r18pa:global}
Assume PG18's local theorem just specified has
passed its separate independent review.
For any preassigned number field $F_{\mathrm{base}}$
and every $D\ge5$, there are a fixed finite
extension $F/F_{\mathrm{base}}$, a rational
prime $q$ unramified in $F$, and a Zariski-dense
set of smooth bidegree-$(D,D)$ divisors
\[
                      X\subset(\mathbf P^1\times\mathbf P^4)_F
 \tag{PA18.6}\label{r18pa:family}
\]
with full ordinary-actual strong G for all
geometrically homologically trivial rational
cycles in all codimensions at every finite
completion.  The assertion holds for all
new cycles after every finite extension
of $F$.  Thus PH16 gives the Beilinson--Bloch
height pairing on the full groups
\[
 \operatorname{CH}^r(X_{F'})_{\mathbf Q,\mathrm{geom\,hom}}
 \times
 \operatorname{CH}^{5-r}(X_{F'})_{\mathbf Q,\mathrm{geom\,hom}}
                       \longrightarrow\mathbf R
 \tag{PA18.7}\label{r18pa:height}
\]
for every finite $F'/F$ and every $1\le r\le4$.
At each place above $q$, \eqref{r18pa:monodromy}
holds.  Every geometric member is not isomorphic
to any smooth projective complete intersection
in projective space.
\end{theorem}

\begin{proof}
\emph{Step 1: fixed smooth reference and one field.}
Choose a bihomogeneous integral $F_0$ of
bidegree $(D,D)$ whose generic divisor is
smooth.  Such sections form a nonempty open
by first-jet surjectivity, or by Bertini in
the Segre embedding.  Its generic divisor
is geometrically integral by the preceding
connectedness calculation.
Let $S_0$ contain $2$ and the finitely many
rational primes of nonsmooth reduction
of this fixed reference section.
At each $p\in S_0$, the smooth-section
open over $\mathbf F_p$ is geometrically
nonempty by the same first-jet incidence
dimension argument.  Choose a point over
some finite extension $\mathbf F_{p^{f_p}}$.
Take a common multiple $N_0\ge2$ of those
degrees and lift, by coefficient CRT,
monic irreducible degree-$N_0$ polynomials
over all those $\mathbf F_p$ to one
monic polynomial over $\mathbf Z$.
Its number field $L$ is unramified and
inert of degree $N_0$ at those primes.
Set $F=F_{\mathrm{base}}L$ and lift the
selected smooth residue sections to forms
$G_p$ over $\mathcal O_L$, hence over
$\mathcal O_F$.  Their reductions are
smooth at every place of $F$ above $p$.
This is exactly the elementary finite-field
and number-field construction in CG18;
no smooth point over the original small
residue field was presumed.

\emph{Step 2: fix the deep condition before the pencil.}
Over $\mathbf Q$ choose $D$ sections
$f_i\in H^0(Y,\mathcal H)$ in geometric
SNC position, with all partial intersections
smooth of expected dimension and no sixfold
intersection, and choose $Q\in V_D$ transverse
to their positive-dimensional strata and
avoiding all their zero-dimensional strata;
one may also require $Q=0$ itself to be
smooth on $Y$.
These are nonempty characteristic-zero
Bertini open conditions in a finite list
of very ample section spaces.
The connectedness calculation above gives
their required geometric integrality.
Clear denominators and put $P=\prod_i f_i$.
Choose an odd rational prime $q$ outside
$S_0$, outside the discriminant of $F$,
and outside the finite set of primes where
the displayed geometric data fail to
spread.  At every $w\mid q$ all the
PG18 local hypotheses on the special
data hold natively and $q$ is a uniformizer.

Put $B=\prod_{p\in S_0}p$, $M=q^2B$,
and choose a bihomogeneous integral $F_1$
subject to
\[
 \begin{aligned}
 F_1&\equiv P+qQ&&\pmod{q^2\mathcal O_F},\\
 F_1&\equiv G_p&&\pmod{p\mathcal O_F}\quad(p\in S_0).
 \end{aligned}
 \tag{PA18.8}\label{r18pa:coefficient-CRT}
\]
This is a full affine coefficient lattice,
hence is Zariski dense over $F$.
The fixed-$F_0$ Lefschetz open is nonempty
on this product ambient: use the first-
and second-jet and separated two-point
first-jet surjectivities supplied by the
Segre embedding above.  The incidence,
invertible-Hessian, finite projection of
the discriminant from $[F_0]$, and
smooth-base arguments in CG18 then apply
with dimension five on this actual $Y$.
They do not require $Y$ to be a complete
intersection.  Choose $F_1$ in that open
and in \eqref{r18pa:coefficient-CRT}.

\emph{Step 3: spread after fixing the pencil and choose parameters.}
Let $S$ be a finite set of prime ideals
over which this chosen pencil fails to
have the smooth total/base and finite
etale separated critical-scheme properties
proved in CG18.  Enlarge $S$ to contain
all primes above $q$ and $S_0$.  Outside
$S$, every geometric member is smooth
or has one ordinary quadratic point.
Exclude primes of residue linear
dependence of the two sections as well.
Let $T$ be the rational primes under $S$.
Choose integers $a$ with
\[
 \begin{gathered}
 a\equiv1\pmod{q^2},\qquad a\equiv1\pmod B,\\
 a\equiv0\pmod p\quad(p\in T\setminus(S_0\cup\{q\})),
 \end{gathered}
 \tag{PA18.9}\label{r18pa:parameter-CRT}
\]
and discard the finitely many singular
generic members.  Set
\[
                           X_a=V(MF_0+aF_1)\subset Y_F.
 \tag{PA18.10}\label{r18pa:equation}
\]
There are infinitely many remaining integers
on each such pencil.

\emph{Step 4: native local equations and complete cycle domains.}
At $w\mid q$ one has the exact equation
\[
                 MF_0+aF_1=P+qQ_a,\qquad
                  Q_a=(MF_0+aF_1-P)/q,\qquad\bar Q_a=\bar Q.
 \tag{PA18.11}\label{r18pa:Qa}
\]
Its coefficients are integral by the two
congruences modulo $q^2$.  PG18 gives
full fourfold G for this actual model,
and \eqref{r18pa:monodromy} gives its
nonzero fourth monodromy.
At $p\in S_0$ the special equation is
the smooth $G_p$ section.  At
$p\in T\setminus(S_0\cup\{q\})$ it
is the smooth $F_0$ section, multiplied
by the unit $M$.  All characteristic-two
places are in these good cases.
At the remaining places, the spread
pencil gives a smooth or one-node
special member in odd characteristic.

Each equation is nonzero on the integral
special ambient fiber.  Since the product
ambient is smooth over the native valuation
ring, this gives a relative Cartier divisor
whose quotient has no uniformizer torsion;
the model is flat and projective.  Its
smooth integral generic fiber makes it
integral.  Its special node members are
geometrically integral by
\eqref{r18pa:connected} and the $R_1,S_2$
argument above.  Thus the frozen ON17
theorem applies at every node place,
with arbitrary finite thickness and no
global complete-intersection requirement.
At good places use the frozen good-
reduction theorem.  This proves every
original local cycle assertion.

After any unramified extension at $w\mid q$,
the PG18 hypotheses hold over the new
native field and its ambient cohomology
is still the same two algebraic divisor
lines.  Apply PG18 afresh to all new
cycles there.  The independently reviewed
RI18 theorem upgrades this to all finite
local extensions, with full-degree
actual trace descent at the last step.
Good reduction and ON17 already apply
afresh to all new cycles after arbitrary
finite extension at every other place.
PH16 and GA17 therefore give
\eqref{r18pa:height} on the full homological
groups, not just on pulled-back cycle
subspaces.

\emph{Step 5: density and the geometric distinction.}
For each allowed $F_1$ infinitely many
parameters $a$ occur.  A coefficient
polynomial vanishing on all resulting
$MF_0+aF_1$ therefore vanishes on the
whole line, and in particular at
$MF_0+F_1$.  The allowed $F_1$ are
dense in a fixed coefficient lattice
and nonempty open, hence in $V_D$.
Translation gives the claimed density
in the full bihomogeneous section space
and its smooth projective open over
this fixed $F$.

For the final geometric assertion use
weak Lefschetz for the smooth ample
divisor $X_a\subset Y$:
\[
 H^2(X_{a,\bar F},\mathbf Q_\ell)
       \simeq H^2(Y_{\bar F},\mathbf Q_\ell)
       \simeq\mathbf Q_\ell(-1)^2.
 \tag{PA18.12}\label{r18pa:b2}
\]
The two classes are the restrictions of
$x,y$, and are independent.  In the
same degree a smooth projective complete
intersection of dimension four has
the one-dimensional hyperplane class
by its weak Lefschetz theorem.  These
geometric Betti numbers are isomorphism
invariants and are unchanged by further
field extension.  The claimed
non-complete-intersection distinction
follows.
\end{proof}

\paragraph{Exact dependencies and unproved extensions.}
PG18 is the only new local theorem activated in
this application; a failed PG18 audit would leave
the product-height conclusion conditional, even
though the arithmetic construction itself remains
valid.  The actual ordered-model and extreme
monodromy calculations are the same local
normal-crossings and point-degree computations
as HP18/HM18, with the positive-stratum
connectedness and the native degree five checked
above.  RI18 supplies arbitrary finite-extension
coverage at the deep places; ON17 handles the
ordinary places; PH16 and GA17 handle complete
heights on that established domain.

The cohomological distinction uses only ordinary
weak Lefschetz (Deligne, \emph{La conjecture de
Weil. II}, Section 4.1.6,
\url{https://www.numdam.org/item/PMIHES_1980__52__137_0/})
and the standard product cohomology calculation.
The complete-intersection comparison is the
same one checked in CA18/CE18.
The positive $N^4$ excludes potential good
reduction and lower-depth dominant witnesses
by frozen NM17; in particular this is not
a change of model for such a generic fiber.
There is no assertion that every divisor in
this linear system belongs to the arithmetic
family, nor that every variety can be brought
into this product ambient.  General G and
general full unconditional heights remain open.
